\chapter{绪论与系统规划}

\section{项目背景}

本项目是一个面向高校教务场景的学生选课成绩管理系统。系统需要支撑学生查询课程、提交选课和退课申请，教师录入成绩并查看成绩分布，管理员维护学期、课程、教学班、学生和教师基础资料。与普通信息展示系统相比，选课系统的核心难点在于数据联系复杂、完整性约束多、并发选课场景明显。课程容量、重复选课、上课时间冲突、先修课程、成绩录入和退课限制都属于典型数据库规则。

项目当前采用 Python + Streamlit 实现界面，使用 psycopg2 连接 openGauss 数据库。数据库部署相关文件位于 \code{opengauss_setup/docker/}，核心建表脚本位于 \code{opengauss_setup/sql/init.sql}。项目没有使用 ORM，也没有使用 Alembic 等迁移框架，数据库模式主要由 SQL 文件维护。

\section{系统规划目标}

根据数据库设计流程，规划阶段首先需要明确系统边界和数据管理目标。本系统的规划目标包括：

\begin{enumerate}
  \item 建立统一的用户账号与角色体系，支持学生、教师、管理员三类用户。
  \item 区分课程定义和具体教学班，支持同一课程在不同学期、不同教师、不同教室开设。
  \item 使用选课记录表达学生和教学班之间的多对多联系，并保存选课状态、选课时间和最终成绩。
  \item 通过结构化上课时间支持时间冲突检查。
  \item 通过主键、外键、唯一约束、CHECK 约束和服务层校验共同保证数据完整性。
  \item 在并发选课场景中使用事务和锁机制降低超选风险。
\end{enumerate}

\section{数据库设计流程说明}

本报告按照数据库设计流程组织章节：需求分析识别角色、功能、数据和规则；概念设计使用 ER 模型抽象实体和联系；逻辑设计把 ER 模型转换为关系模式；规范化设计分析函数依赖和范式；完整性设计讨论主键、外键、CHECK、UNIQUE 和应用层校验；事务与锁章节分析并发选课一致性；物理设计讨论数据类型、索引和性能；运行维护章节说明初始化、部署、安全和扩展。

\section{本章课程知识点}

本章对应数据库系统设计流程中的规划阶段。规划阶段强调数据库系统建设的必要性、系统边界、软硬件环境、数据范围和设计目标。本项目的数据中心特征明显：无论学生选课、教师录入成绩还是管理员维护教学班，本质上都是围绕数据库中实体、联系和约束展开的操作。
